\section{Связанные исследования}\label{sec:related}

\paragraph{Выбор разбиения по ограниченной внешней информации.}
В полуобучаемой кластеризации известные метки или парные ограничения могут использоваться как при построении кластеров, так и при выборе готового результата. Pourrajabi и соавторы~\citep{pourrajabi2014} рассматривают выбор модели и организацию проверки с учётом зависимости ограничений. COBS~\citep{cobs2017} использует ограничения для выбора между результатами разных алгоритмов и настроек, полученными без учителя. Авторы показывают экспериментально, что такой выбор способен превосходить сравниваемые полуобучаемые методы. Это непосредственный предшественник постановки с фиксированным пулом: внешняя информация направляет выбор, не изменяя самих кандидатов. В настоящей работе внешним сигналом служит ARI на подмножестве объектов, а предметом анализа --- последствия допуска альтернатив с учётом изменчивости этого сигнала.

\paragraph{Сочетание критериев качества.}
Akhanli и Hennig~\citep{akhanli2020} рассматривают внутренние индексы, характеризующие различные свойства кластеров, и калибруют их относительно случайных разбиений тех же данных перед агрегацией. Предложенные агрегированные индексы сравниваются на синтетических и реальных данных. Этот подход показывает, как сделать сопоставимыми оценки разных шкал, но не оценивает неопределённость внешнего критерия по малой разметке. Более близкий по источникам информации подход представлен в работе Castro Gertrudes и соавторов~\citep{gertrudes2019}: при извлечении кластеров из плотностной иерархии сочетаются структурная устойчивость и внешняя оценка на основе меток, включая неполное представление категорий. Следовательно, объединение внутренней и внешней информации и проблема отсутствующих классов имеют прямых предшественников. Здесь рассматривается выбор среди готовых разбиений разных алгоритмов и отдельно проверяется роль бутстрепного допуска.

\paragraph{Зависимость внутренних индексов от условий задачи.}
Benchmark-исследование Simpson и соавторов~\citep{simpson2026benchmark} сравнивает 26 внутренних индексов с использованием нескольких взаимодополняющих протоколов. Сопоставление выбранного лучшего разбиения и ранжирования кандидатов позволяет анализировать разные аспекты поведения индекса. В работе об instance space~\citep{simpson2026instance} девять индексов исследуются на 18\,351 синтетическом наборе и результатах восьми алгоритмов. Выявленные области успешной работы индексов мотивируют изучение условий их применимости вместо единого общего рейтинга. Для нашей задачи дополнительно существенны доступность внешних меток и величина потери после окончательного выбора: сходство ранжирований само по себе не определяет последствия конкретного решения.

\paragraph{Неопределённость после выбора.}
Efron~\citep{efron2014} рассматривает оценивание точности после выбора модели. Эта работа служит методологическим основанием учитывать влияние отбора на статистический вывод, но не обосновывает используемый здесь порог. Бутстреп относительно лидера, выбранного на той же размеченной выборке, характеризует условную изменчивость сравнения. Он не гарантирует правильности выбора при смещённой разметке и не восстанавливает отсутствующие классы.

Перечисленные направления определяют контекст исследования: выбор по ограниченной внешней информации, сочетание разных характеристик качества и анализ условий полезности индексов уже изучаются. Настоящая работа сосредоточена на проверке адаптивного допуска и диагностике его вмешательств в выбор внешнего лидера. Обзор имеет тематический характер и не является доказательством отсутствия аналогичного правила во всей литературе.
